\section{Compito due: raggiungibilità e distanza libera}

Il secondo compito si concentra sull'analisi topologica dello spazio locale a partire da una cella di origine $O$. Vengono introdotti i concetti geometrici di distanza libera ($d_{\text{lib}}$), cammino libero di tipo 1 e tipo 2, contesto, complemento e frontiera della chiusura.

\subsection{Distanza libera ($d_{\text{lib}}$) e formulazione matematica}

La distanza libera ($d_{\text{lib}}$) rappresenta il costo minimo del percorso ottimale tra due celle in una griglia con adiacenza a otto direzioni, assumendo l'assoluta assenza di ostacoli fisici interposti. Siano $O$ e $D$ le celle di origine e destinazione, e siano $\Delta x = |O.x - D.x|$ e $\Delta y = |O.y - D.y|$ le differenze assolute delle loro coordinate cartesiane sulla griglia.

Definiamo $\Delta_{\text{min}} = \min(\Delta x, \Delta y)$ come il numero massimo di passi diagonali puri eseguibili e $\Delta_{\text{max}} = \max(\Delta x, \Delta y)$ come la dimensione dominante dello spostamento. La formula matematica esatta per il calcolo di $d_{\text{lib}}(O,D)$ è:

\begin{equation}
d_{\text{lib}}(O,D) = \sqrt{2}\Delta_{\text{min}} + (\Delta_{\text{max}} - \Delta_{\text{min}})
\label{eq:dlib}
\end{equation}

Questa formula di costo $O(1)$ fornisce il limite inferiore richiesto dalla potatura geometrica descritta di seguito.

\subsection{Contesto e complemento}

La raggiungibilità locale tramite cammini liberi definisce due insiemi disgiunti di celle all'interno dello spazio libero:
\begin{itemize}
    \item \textbf{Contesto}: l'insieme delle celle raggiungibili da $O$ tramite un cammino libero di \textbf{tipo 1}, definito dalla sequenza prima mosse diagonali oblique, poi mosse cardinali residue, oppure da cammini composti esclusivamente da un tipo di mossa.
    \item \textbf{Complemento}: l'insieme delle celle raggiungibili da $O$ \textbf{solo} tramite un cammino libero di \textbf{tipo 2}, definito dalla sequenza inversa prima mosse cardinali, poi mosse diagonali residue. Le celle che ammettono entrambi i tipi di cammino appartengono per convenzione al solo contesto.
\end{itemize}
L'unione di contesto e complemento forma la \textbf{chiusura} di $O$.

\subsection{Proiezione di raggi geometrici ($O(\max(R,C)^2)$)}

Per determinare gli insiemi contesto e complemento, il programma implementa la \textbf{proiezione di raggi}. L'analisi controlla le linee di raggiungibilità procedendo a raggiera a partire dall'origine, riducendo la complessità da quadratica ($O(R^2 \cdot C^2)$ per una ricerca cella per cella) a $O(\max(R,C)^2)$.

L'algoritmo per il \textbf{contesto (tipo 1)} opera come segue:
\begin{enumerate}
    \item \textbf{Proiezione cardinale}: vengono tracciati i 4 raggi cardinali puri (Nord, Sud, Est, Ovest) a partire dall'origine $O$, arrestandosi non appena si incontra una cella ostacolo.
    \item \textbf{Proiezione diagonale ed espansione a gomito}: per ciascuno dei 4 quadranti del piano cartesiano locale:
    \begin{itemize}
        \item si proietta il raggio diagonale puro (NE, NO, SE, SO);
        \item per ogni cella attraversabile raggiunta lungo questo raggio diagonale, si avvia un'espansione cardinale rettilinea parallela agli assi (espansione a gomito), che procede verso l'esterno fino a incontrare un ostacolo o il bordo della griglia.
    \end{itemize}
\end{enumerate}

L'algoritmo per il \textbf{complemento (tipo 2)} inverte simmetricamente le fasi di proiezione: per ciascuno dei 4 quadranti si proiettano prima i raggi cardinali e, da ogni cella libera raggiunta su di essi, si espandono raggi in diagonale pura. Al termine si escludono le celle già presenti nel contesto.

Questo accorgimento analizza ogni cella limitando le operazioni di lettura sulla matrice.

Di seguito si riporta lo pseudocodice dei due algoritmi. La scrittura $P + c$ indica la cella ottenuta spostando $P$ di un passo lungo la direzione $c$; una cella si dice \emph{libera} se è interna alla griglia e non appartiene a $osta$.

\begin{algorithm}[H]
\caption{Calcolo del contesto (celle raggiungibili con cammini liberi di tipo 1)}
\label{alg:calcola_contesto}
\begin{algorithmic}[1]
\Procedure{CalcolaContesto}{$O, dim, osta$}
    \State $contesto \leftarrow \{O\}$
    \For{\textbf{each} direzione cardinale $c \in \{\text{Nord}, \text{Sud}, \text{Est}, \text{Ovest}\}$}
        \State $P \leftarrow O + c$
        \While{$P$ è libera}
            \State $contesto \leftarrow contesto \cup \{P\}$
            \State $P \leftarrow P + c$
        \EndWhile
    \EndFor
    \For{\textbf{each} direzione diagonale $d \in \{\text{NE}, \text{NO}, \text{SE}, \text{SO}\}$}
        \State $P \leftarrow O + d$
        \While{$P$ è libera}
            \State $contesto \leftarrow contesto \cup \{P\}$
            \State \Comment{Espansione cardinale orizzontale e verticale a partire da $P$}
            \For{\textbf{each} $c \in \{\text{componente orizzontale di } d, \text{componente verticale di } d\}$}
                \State $Q \leftarrow P + c$
                \While{$Q$ è libera}
                    \State $contesto \leftarrow contesto \cup \{Q\}$
                    \State $Q \leftarrow Q + c$
                \EndWhile
            \EndFor
            \State $P \leftarrow P + d$
        \EndWhile
    \EndFor
    \State \Return $contesto$
\EndProcedure
\end{algorithmic}
\end{algorithm}

\begin{algorithm}[H]
\caption{Calcolo del complemento (celle raggiungibili solo con cammini liberi di tipo 2)}
\label{alg:calcola_complemento}
\begin{algorithmic}[1]
\Procedure{CalcolaComplemento}{$O, dim, osta, contesto$}
    \State $complemento \leftarrow \emptyset$
    \For{\textbf{each} direzione diagonale $d \in \{\text{NE}, \text{NO}, \text{SE}, \text{SO}\}$}
        \For{\textbf{each} $c \in \{\text{componente orizzontale di } d, \text{componente verticale di } d\}$}
            \State $P \leftarrow O + c$
            \While{$P$ è libera}
                \State \Comment{Espansione diagonale a partire dalla cella $P$ del raggio cardinale}
                \State $Q \leftarrow P + d$
                \While{$Q$ è libera}
                    \If{$Q \notin contesto$}
                        \State $complemento \leftarrow complemento \cup \{Q\}$
                    \EndIf
                    \State $Q \leftarrow Q + d$
                \EndWhile
                \State $P \leftarrow P + c$
            \EndWhile
        \EndFor
    \EndFor
    \State \Return $complemento$
\EndProcedure
\end{algorithmic}
\end{algorithm}

\subsection{Frontiera della chiusura}

La \textbf{frontiera} della chiusura di $O$ è definita formalmente come il sottoinsieme delle celle della chiusura che confinano (secondo l'intorno a otto direzioni) con almeno una cella attraversabile che si trova all'esterno della chiusura stessa:

\begin{equation}
\text{Frontiera}(O) = \{ c \in \text{Chiusura}(O) \mid \exists v \in \text{Vicini}(c) \text{ t.c. } \text{Stato}(v) = 0 \land v \notin \text{Chiusura}(O) \}
\label{eq:frontiera}
\end{equation}

Il codice estrae la frontiera limitando l'ispezione al vicinato delle sole celle incluse nella chiusura. Ad ogni cella di frontiera viene associato il tipo del cammino libero da cui proviene ($1$ per il contesto, $2$ per il complemento).
